\phantomsection\addcontentsline{toc}{section}{Sustainable Buildings}
\ECchapter{29}{Sustainable Buildings}{Designing low-carbon and comfortable spaces}{ECOrange}

\ECObjectives{
\begin{itemize}
\item Explain the main physical and engineering principles.
\item Interpret measurements, models and performance trade-offs.
\item Identify safety, environmental and validation constraints.
\end{itemize}}

\begin{multicols}{2}
\begin{engineeringnews}
Buildings consume energy during construction and throughout decades of operation. Sustainable design therefore combines passive architecture, efficient equipment, low-carbon materials, renewable energy and monitoring of real performance.
\end{engineeringnews}

\begin{ecarticle}
A sustainable building must provide comfort, safety and useful space while reducing environmental impacts over its entire lifecycle. Operational energy is used for heating, cooling, ventilation, lighting and equipment. Embodied carbon is associated with extracting materials, manufacturing products, transport, construction, maintenance and demolition. As grids become cleaner and insulation improves, embodied impacts represent a larger share of the total.

The first design strategy is to reduce demand. Building orientation, window area, shading and thermal mass influence solar gains. Insulation and airtightness reduce heat transfer, while controlled ventilation maintains air quality. External shading can block summer sun without preventing winter daylight. These passive decisions are usually more robust than adding complex equipment later.

Heating and cooling systems must then match the reduced demand. Heat pumps can transfer several units of heat for one unit of electricity, but their performance depends on outdoor temperature and supply temperature. Efficient fans, pumps and controls reduce auxiliary consumption. Sensors can adjust operation according to occupancy, carbon-dioxide concentration and weather forecasts.

Material selection requires lifecycle assessment. Concrete offers durability and thermal mass but cement production emits carbon dioxide. Timber stores biogenic carbon and can reduce mass, yet sourcing, fire safety, moisture and end-of-life must be considered. Reuse of structural elements can avoid both waste and new production.

A design model is only a prediction. The performance gap between calculated and measured consumption may result from construction defects, incorrect controls or occupant behaviour. Commissioning, metering and feedback are therefore part of sustainable engineering. A successful building is not merely certified at handover; it continues to perform well during use.
\end{ecarticle}

\begin{tcolorbox}[title=\sffamily\bfseries Did You Know?,colback=yellow!10,colframe=ECOrange]
A high-performance prototype is not sufficient: engineers must prove repeatability, safety and useful performance under realistic operating conditions.
\end{tcolorbox}

\begin{engineeringfocus}
\textbf{Heat and energy flows}
\begin{center}
\begin{tikzpicture}[font=\scriptsize,>=Latex]
\draw[thick,fill=ECOrange!8] (1,0) rectangle (5,2.8);
\draw[fill=ECLightBlue] (2.7,0) rectangle (3.3,1.2);
\node at (3,1.8) {indoor space};
\draw[->,thick,ECRed] (0,2.2)--node[above]{solar gain}(1,2.2);
\draw[->,thick,ECBlue] (5,2.0)--node[above]{heat loss}(6.7,2.0);
\draw[->,thick,ECGreen] (3,-.6)--node[right]{heat pump}(3,0);
\draw[->,thick,ECPurple] (4.2,3.5)--node[right]{ventilation}(4.2,2.8);
\end{tikzpicture}
\end{center}
\end{engineeringfocus}

\begin{keyfigures}
\begin{tabularx}{\linewidth}{@{}Xr@{}}
First priority & reduce demand\\
Envelope metric & U-value\\
Air quality proxy & CO$_2$\\
Operational check & metering\\
Material tool & lifecycle assessment\\
\end{tabularx}
\end{keyfigures}

\begin{tcolorbox}[title=\sffamily\bfseries Illustrative annual energy demand,colback=white,colframe=ECOrange]
\begin{center}
\begin{tikzpicture}
\begin{axis}[ybar stacked,bar width=13pt,width=\linewidth,height=50mm,ylabel={Energy (kWh/m$^2$/year)},symbolic x coords={Conventional,Efficient,Passive},xtick=data,legend style={font=\scriptsize,at={(0.5,-0.22)},anchor=north,legend columns=3},ticklabel style={font=\scriptsize}]
\addplot table[x=design,y=heating,col sep=comma]{data/EC29/energy.csv};
\addplot table[x=design,y=cooling,col sep=comma]{data/EC29/energy.csv};
\addplot table[x=design,y=lighting,col sep=comma]{data/EC29/energy.csv};
\legend{Heating,Cooling,Lighting}
\end{axis}
\end{tikzpicture}
\end{center}
\end{tcolorbox}

\begin{vocabulary}
\begin{tabularx}{\linewidth}{@{}lX@{}}
embodied carbon & carbone incorporé\\
operational energy & énergie d’exploitation\\
airtightness & étanchéité à l’air\\
thermal bridge & pont thermique\\
shading & protection solaire\\
heat pump & pompe à chaleur\\
commissioning & mise au point\\
metering & comptage\\
performance gap & écart de performance\\
lifecycle assessment & analyse du cycle de vie\\
\end{tabularx}
\end{vocabulary}

\begin{questions}
\item What is the difference between operational and embodied carbon?
\item Why should demand reduction precede equipment selection?
\item How can external shading improve summer comfort?
\item Why can actual consumption differ from simulation?
\item Give two reasons for monitoring indoor carbon dioxide.
\end{questions}

\begin{engineeringchallenge}
Design a low-energy classroom for a temperate climate. Propose envelope targets, shading, ventilation, heating, sensors and a method for checking actual performance after one year.
\end{engineeringchallenge}

\begin{speaking}
Should building regulations limit embodied carbon as strictly as operational energy?
\end{speaking}
\end{multicols}

\subsection*{Sources}
\ECsource{International Energy Agency -- Buildings}{https://www.iea.org/energy-system/buildings}
\ECsource{UN Environment Programme -- Buildings and construction}{https://www.unep.org/topics/cities/buildings-and-construction}
